\documentclass[10pt,a4paper]{amsart}
\usepackage[margin=19mm]{geometry}
\usepackage{amsmath,amssymb,mathtools}
\usepackage[scr=rsfs,cal=euler]{mathalfa}
\usepackage{xfrac}
\usepackage[unicode,hidelinks,bookmarks=false]{hyperref}
\newcommand{\ie}{i.e.\ }
\newcommand{\cf}{cf.\ }
\newcommand{\resp}{respectively\ }
\newcommand{\Subsection}{\subsection}
\providecommand{\nobreakdash}{\nobreak\hbox{-}}
\setlength{\emergencystretch}{2em}
\multlinegap0pt
\usepackage{amssymb}
\usepackage{tikz}
\newtheorem{definition}{Definition}
\newtheorem{theorem}{Theorem}
\newtheorem{lemma}{Lemma}
\title{Maximal entropy dissipation numerical scheme for conservation law systems}
\author{Marko Nedeljkov}
\date{Source: arXiv:2608.15716v1 (CC BY 4.0)}
\begin{document}
\maketitle

\noindent\emph{Source and licence.} Maximal entropy dissipation numerical scheme for conservation law systems. Version arXiv:2608.15716v1 (CC BY 4.0), distributed by arXiv under the Creative Commons Attribution 4.0 International licence, http://creativecommons.org/licenses/by/4.0/, as recorded in the arXiv licence metadata for this exact version and shown on the abstract page https://arxiv.org/abs/2608.15716v1. Full licence notice: \texttt{licenses/arxiv-2608.15716-CC-BY-4.0.txt}.\par\smallskip

\noindent\textit{Editorial scope.} Counted unit: the article's theorem t-2 (source0.tex lines 661--666). The article's complete original proof of it is delivered with it (source0.tex lines 668--732, one \texttt{proof} environment of 65 lines). No mathematics is written, repaired or re-derived: the statement and the proof are the article's own lines, in the article's own notation, and no retained line is substituted or re-flowed.  Retained uncounted as the source-local context the proof needs: lines 141--146; lines 192--207; lines 228--231; lines 313--341; lines 396--400; lines 419--422; lines 447--459; lines 466--492; lines 618--623.  Nothing was omitted from any retained span, so the package carries no omission record.  No retained line contains a citation, so the wrapper carries no bibliography and no bibliography entry.  No retained line contains a figure environment, an image-inclusion command or drawing code, so no image is delivered and the package declares no asset dependency.  Editorial formatting only, in addition: an AMS \texttt{amsart} preamble that restates the article's own declarations and macros used by the retained lines, namely the environments \texttt{definition}, \texttt{theorem}, \texttt{lemma} on separate counters, together with the standard packages those lines need.  The retained environments print in this document as Definition 1, Theorem 1, Lemma 1 in order of appearance, whereas the article prints the same items with the numbering of its own full document.  The complete native source of the article as distributed in the arXiv e-print for this exact version is bundled unchanged as \texttt{sources/207775/source0.tex}.
\begin{equation} \label{claws}
% U_{t}+\nabla \cdot F(U) = \Pi(U) \text{ in } \mathbb{R}_{+}^{d}
%:=\mathbb{R} ^{d}\times  (0,\infty ), 
U_{t}+\nabla \cdot F(U) = 0 \text{ in } \mathbb{R}_{+}^{m}:=\mathbb{R} ^{m}
\times  (0,\infty ), \;  U(x,0)=U_{0}(x), 
\end{equation}

\begin{definition} \label{def-dafermos}
A weak solution $U$ to system (\ref{claws}) satisfies the energy
admissibility condition if for any other weak solution
$\bar{U}$ such that for some $\tau\in [0,T)$ we have
$U(\cdot,t)=\bar{U}(\cdot,t)$, $t\leq \tau$ and 
there exists a sequence $\{\tau_{i} \}_{i \in \mathbb{N}}$, $\tau_{i}>\tau$,
$\tau_{i}\to \tau$, as $n\to \infty$ such that
\[
\begin{split}
& \mathcal{E}(\bar{U})(\tau +)\geq \mathcal{E}(U)(\tau +), \\ 
& \text{where }
\mathcal{E}(U)(\tau +)=\mathop{\rm esslim}_{t \to \tau, t>\tau}\mathcal{E}(U)(t)
\text{ in the distributional sense}.
\end{split}
\]
\end{definition}

\begin{equation} \label{shceme}
U_{i}^{j+1}=U_{i}^{j}-\dfrac{\Delta t}{|\Omega _{i}| }
\sum_{k\in K_{i}}\int_{S_{i,k}}h_{F\cdot \nu_{i,k}}(U_{i}^{j},U_{k}^{j})\, dS,
\end{equation}

\begin{definition} \label{defscheme}
For a given initial data $U_{0\Delta}$ we construct approximate 
(i.e.\ numerically obtained) solution 
$U_{\Delta}$ in $D_{T}$ by (\ref{shceme}) where
\begin{equation} \label{nflux}
h_{F\cdot \nu_{i,k}}(U_{i}^{j},U_{k}^{j})
=(\theta_{i,k}^{j} \cdot \underline{F}_{i,k}^{j}+(1-\theta_{i,k}^{j})\cdot
\overline{F}_{i,k}^{j})\cdot \nu_{i,k}, 
\end{equation}
where 
\[
\begin{split}
\underline{F}_{i,k}^{j}&:=\min\{F(U):\; U\in [\min\{ U_{i}^{j},
U_{k}^{j}\},\max\{U_{i}^{j}, U_{k}^{j}\}]\}\\
\overline{F}_{i,k}^{j}&:=\max\{F(U):\; U\in [\min\{ U_{i}^{j},
U_{k}^{j}\},\max\{U_{i}^{j}, U_{k}^{j}\}]\},
\\ & \theta_{i,k}^{j}=\theta_{k,i}^{j}.
\end{split}
\] 	
For each $j\in \mathbb{N}_{0}$, $i,k\in I$, 
$\theta_{i,k}^{j}$ is $n$-dimensional vector with values in $[0,1]$, 
$1$ denotes the $n$-dimensional vector with each component equals $1$, 
and ``$\cdot$'' is the usual scalar product. 
The functions $\min$ and $\max$ are used component-wise.
	
The elements $\{\theta_{i,k}^{j}\}_{i,k\in I}\in [0,1]^{n}$ are chosen such that 
\[\eta^{j+1}(U_{\Delta}):=\int_{\Omega_{T}}\eta(U_{\Delta}(x,t_{j+1}))\,dx\] 
is minimal for each time step $j$.
\end{definition}

\begin{theorem} \label{lw}
Let $U_{\Delta}:D_{T} \to B_{T}$, assume $B_{T}$ is bounded and 
$\eta\in \mathcal{C}^{2}(B_{T})$.
If $U_{\Delta} \to U$ boundedly almost everywhere on $D_{T}$ as $|\Delta|\to 0$, then the limit function $U$ is a weak solution of (\ref{claws}).
\end{theorem}

\begin{equation} \label{if}
\int ^{x_{2}}_{x_{1}}U( x,t_{2}) -U( x,t_{1}) \, dx+\int ^{t_{2}}_{t_{1}}
F( U( x_{2},t) ) -F( U( x_{1},t) ) \, dt=0,
\end{equation}

\begin{equation} \label{1dscheme}
\begin{split}
U_{i}^{j+1} =  & U_{i}^{j}+\Big( \theta_{i}^{j}\underline{F}_{i-1/2}^{j}
+(1-\theta_{i}^{j})\overline{F}_{i-1/2}^{j}) \\
& -\big(\theta_{i+1}^{j}\underline{F}_{i+1/2}^{j}
+(1-\theta_{i+1}^{j})\overline{F}_{i+1/2}^{j}\big) \Big) \frac{\Delta t}{\Delta x}\\
& \underline{F}_{i-1/2}^{j}=\min\{F(U):\; 
U\in [\min\{U_{i-1},U_{i}\},\max\{U_{i-1},U_{i}\}]\} \\
& \overline{F}_{i-1/2}^{j}=\max\{F(U):\; 
U\in [\min\{U_{i-1},U_{i}\},\max\{U_{i-1},U_{i}\}]\},\\ 
& i\in I,\; j\in \mathbb{N}.
\end{split}
\end{equation}

\begin{lemma} \label{lemma3.1}
Let $\widetilde{U}$ be the piecewise constant approximation of the classical weak solution
$\widehat{U}$ defined on the equidistant mesh 
$\{ (x_{i-1/2},t_{j})\}_{i\in I, j\in \mathbb{N}}$, $\Delta x=x_{i+1/2}-x_{i-1/2}$, 
$\Delta t=t_{j+1}-t_{j}$ using (\ref{if}):
\[
\begin{split}
& \widetilde{U}(x,t)=\widehat{U}_{i}^{j},\; (x,t) \in (x_{i-1/2},x_{i+1/2})\times [t_{j},t_{j+1}) \\
& F(\widetilde{U}(x_{i+1/2},t))=\widehat{F}_{i+1/2,j},\; t\in [t_{j},t_{j+1}). 
\end{split} 
\]
where 
\begin{equation}\label{kappa} 
\begin{split}
& \widehat{U}_{i}^{j}:= \frac{1}{x_{i+1}-x_{i}}\int_{x_{i-1/2}}^{x_{i+1/2}} \widehat{U}(x,t_{j}) \, dx \\
& \widehat{F}_{i+1/2}^{j}:=\frac{1}{t_{j+1}-t_{j}}\int ^{t_{j+1}}_{t_{j}}F(\widehat{U}(x_{i+1/2},t)) \, dt.
\end{split}
\end{equation}

Then for every $j\in \mathbb{N}_{0}$ 
there exists a fine enough mesh and a FVM approximate solution 
$\breve{U}$ defined at the same mesh and satisfying (\ref{1dscheme}) such that 
\begin{equation}\label{s}
\sum_{i\in I}|\eta(\breve{U}_{i}^{j+1})-\eta(\widetilde{U}_{i}^{j+1})|\Delta x
\end{equation}
is arbitrary small.
\end{lemma}

\begin{equation} \label{hatlinf}
\sum_{i\in I\setminus I_{\rm nd}^{j}\setminus I_{\rm nd}^{j+1}
\setminus I_{\rm m}^{j}}
\| \widehat{U}(x,t_{j}) - \widehat{U}_{i}^{j}\|_{L^{\infty}(D_{T})}{\Delta x} 
< \varepsilon.
\end{equation}

\begin{theorem} \label{t-2}
There is a non-trivial time interval $[0,T_{1}]$ where the limiting solution 
$U=\lim_{\Delta \to 0}U_{\Delta}$ from Theorem \ref{lw} is the one 
satisfying Definition \ref{def-dafermos} among all classical solutions with values 
in the bounded set $B_{T}$.
\end{theorem}

\begin{proof}

Suppose that $U$ is an MES solution to (\ref{claws}), that is
a weak solution obtained as the limit of 
the approximations $U_{\Delta}$ that exist because of Theorem \ref{lw}. 
Suppose $\widehat{U}$ is a classical solution to the system (\ref{claws}) 
that violates the conditions of the Definition \ref{def-dafermos} 
after a time $t=T_{0}$. That is, $\mathcal{E}(U)(T_{0})
=\mathcal{E}(\widehat{U})(T_{0})$ and for every sequence
$(T_{k})_{k\in \mathbb{N}}$  converging to $T_{0}$, $T_{k}>T_{0}$,
\[
\begin{split}
& \int_{-L}^{L} \eta(U(x,t)) dx \leq \int_{-L}^{L} \eta(\widehat{U}(x,t))dx,
\; t\leq T_{0} \text{ and }\\ 
& \int_{-L}^{L} \eta(U(x,T_{k}))dx - \int_{-L}^{L} 
\eta(\widehat{U}(x,T_{k}))dx >\mu>0.
\end{split}
\]
Like in Lemma \ref{lemma3.1}, using the integral form (\ref{if}), we approximate $\hat{U}$ by taking
\begin{equation} \label{hatscheme}
\widehat{U}_{i}^{j+1} = \widehat{U}_{i}^{j} 
-\dfrac{\Delta t}{\Delta x}(\widehat{F}_{i+1/2}^{j}-\widehat{F}_{i-1/2}^{j}).
\end{equation}

Choose $T_{k}$ such that $\Delta t=T_{k}-T_{0}<\mu$, 
One can make the approximation arbitrary fine, so we may take $\Delta t$ small enough 
such that 
$$
\Big| \int_{-L}^{L}\eta(\hat{U}(x,T_{k}))dx-\sum_{i\in I}\eta(\hat{U}_{i}^{j})\Delta x
\Big| < \varepsilon, \; t_{j}=T_{k}.
$$
when $|\Delta|\to 0$, which determines the solution $U$.
By (\ref{hatlinf}), for 
$\varepsilon$ small enough, we may change the definition of $\widehat{U}_{i}^{j}$,
$i\in I$, $j=1,2, \ldots$ for the cells with indices belonging to 
$I \setminus I_{\rm nd}^{j}\cup I_{\rm nd}^{j+1}\cup I_{\rm m}^{j}$ and obtain the approximation
$\{\breve{U}_{i}^{j}, \; i\in I, \; j=1,2, \ldots\}$ so that
the flux function in (\ref{hatscheme}) satisfies the condition (\ref{nflux}).
That is, $(\breve{U}_{i}^{j})_{i\in I}$ is obtained by the scheme from 
Definition \ref{defscheme} for some choice of $\theta_{i}^{j}\in [0,1]^{m}$, 
$i\in I$, $j\geq 0$. So, for the same data at $t=T_{k}=$, 
$\eta^{j+1}(U_{\Delta})\leq \eta^{j+1}(\hat{U})$, where we have used the notation
$\eta^{j}(V):=\sum_{i\in I} \eta(V_{i}^{j})$ for a piecewise constant
approximation $V_{i}^{j}$, $i\in I$, $j=1,2, \ldots$ of the function $V$.

At the same time, using the differentiability of $\eta$, 
boundedness od $\widehat{U}$, and (\ref{hatlinf}), we have the following
inequality
\[
\begin{split}
& \Big| \int_{\Omega_{T}} \Big(\eta(\widehat{U}(\cdot,t_{j+1}))
- \eta^{j+1}(\widetilde{U})\Big) dx \Big| 
= \Big| \int_{\Omega_{T}}\eta(\widehat{U}(\cdot,t_{j+1}))dx 
- \sum_{i\in I}\eta(\widehat{U}_{i}^{j+1})\Delta x \Big| \\
\leq & \|\nabla \eta (\widehat{U})\|_{L^{\infty}(D_{T})}
\sum_{i\in I} \|\widehat{U}(x,t_{j})-\widehat{U}_{i}^{j}\|_{L^{\infty}(D_{T})} \Delta x
< L_{B}\varepsilon,
\end{split}
\]  
for $\Delta t<\mu$, where $L_{B}:=\|\nabla \eta (\widehat{U})\|_{L^{\infty}(D_{T})}
<\infty$.

That contradicts the fact that $\int_{-L}^{L} \eta(U(\cdot,T_{k})) 
- \int_{-L}^{L} \eta(\widehat{U}(\cdot,T_{k}))>\mu$ if $U_{\Delta}=\widetilde{U}$ in $t=T_{0}$.
\end{proof}

\end{document}
